\documentclass[11pt]{article}
\usepackage[margin=1in]{geometry}
\usepackage{amsmath,amssymb}
\usepackage{array}
\usepackage{enumitem}
\usepackage{hyperref}
\usepackage[T1]{fontenc}

\setlength{\parindent}{0pt}
\setlength{\parskip}{0.6em}
\newcommand{\unchecked}{\(\square\)}

\title{Quantum Circuit Design: Weekly Study Notes}
\author{Alexandra Bouillon}
\date{October 12--18, 2026}

\begin{document}
\maketitle

\section{Goal}

Learn to design a circuit that communicates five things clearly:

\begin{enumerate}
    \item Which node owns each qubit.
    \item Which qubits are already entangled.
    \item Which gates happen locally.
    \item Which measurements produce classical bits.
    \item Which later gates are controlled by those classical results.
\end{enumerate}

\section{Design Rules}

\subsection{Label wires by owner and role}

Prefer:

\begin{verbatim}
A.data       Alice's unknown input state
A.link       Alice's half of a Bell pair
S.left       Switch half linked to Alice
S.right      Switch half linked to Bob
B.link       Bob's half of a Bell pair
\end{verbatim}

Avoid labels such as \(q_0,q_1,q_2,q_3\) alone once multiple nodes are involved. Numbered labels hide ownership.

\subsection{Separate the three planes}

\begin{itemize}
    \item \textbf{Quantum state:} qubit wires, gates, memories, and measurements.
    \item \textbf{Quantum channel:} physical process that establishes a Bell link between neighboring nodes.
    \item \textbf{Classical control:} heralding, link-state messages, BSM outcomes, and correction commands.
\end{itemize}

A clear diagram should not make classical messages look like transmitted qubits.

\subsection{Design in blocks}

\begin{center}
\fbox{Prepare links}
\(\rightarrow\)
\fbox{Herald/store}
\(\rightarrow\)
\fbox{BSM}
\(\rightarrow\)
\fbox{Send result}
\(\rightarrow\)
\fbox{Correct}
\(\rightarrow\)
\fbox{Use ebit}
\end{center}

This matches the network sequence:

\begin{itemize}
    \item External links are established first.
    \item Their success states are shared classically.
    \item The switch performs an internal swap.
    \item The endpoints interpret or correct the resulting Bell state.
\end{itemize}

\section{Bell State Measurement Block}

A standard logical BSM circuit for two qubits is:

\begin{verbatim}
S.left:  ----o----H----M---- c0
             |
S.right: ----X---------M---- c1
\end{verbatim}

The CNOT and Hadamard rotate the Bell basis into the computational basis. Measurement then identifies one of four Bell-state outcomes using two classical bits.

\subsection{Correction lookup}

One common convention is:

\begin{center}
\begin{tabular}{|c|c|}
\hline
\textbf{BSM result} & \textbf{Endpoint correction} \\
\hline
\(\Phi^+\) & \(I\) \\
\(\Phi^-\) & \(Z\) \\
\(\Psi^+\) & \(X\) \\
\(\Psi^-\) & \(XZ\) \\
\hline
\end{tabular}
\end{center}

The exact bit-to-gate mapping depends on wire ordering and circuit convention. State the convention instead of memorizing an unlabeled table.

\section{One-Switch Design Template}

\subsection{Inputs}

\begin{itemize}
    \item Bell link between \texttt{A.link} and \texttt{S.left}.
    \item Bell link between \texttt{S.right} and \texttt{B.link}.
\end{itemize}

\subsection{Local switch operation}

\begin{enumerate}
    \item Confirm both links succeeded.
    \item Apply CNOT from \texttt{S.left} to \texttt{S.right}.
    \item Apply \(H\) to \texttt{S.left}.
    \item Measure both switch qubits.
    \item Send the two-bit result over a classical channel.
\end{enumerate}

\subsection{Outputs}

\begin{itemize}
    \item \texttt{S.left} and \texttt{S.right} are measured and freed.
    \item \texttt{A.link} and \texttt{B.link} are entangled, up to a known Pauli frame or correction.
    \item The switch has not measured either endpoint qubit.
\end{itemize}

\section{Circuit-Design Questions}

Ask these before accepting a diagram:

\begin{itemize}[leftmargin=*]
    \item[\unchecked] Are all qubits initialized or supplied by a named link?
    \item[\unchecked] Does every two-qubit gate act on qubits at the same physical node?
    \item[\unchecked] If a gate appears to cross nodes, is it actually a link-generation abstraction?
    \item[\unchecked] Are successful links heralded before the BSM starts?
    \item[\unchecked] Are memory and coherence assumptions visible?
    \item[\unchecked] Are measurements clearly destructive?
    \item[\unchecked] Are classical outcomes drawn as classical wires or messages?
    \item[\unchecked] Is the correction convention stated?
    \item[\unchecked] Are consumed qubits distinguished from surviving endpoint qubits?
    \item[\unchecked] Is the desired output state written explicitly?
\end{itemize}

\section{Common Design Mistakes}

\begin{itemize}
    \item \textbf{Drawing a CNOT directly across distant nodes.} Ordinary circuit notation can hide that a nonlocal operation requires an entanglement-assisted protocol or physical interaction.
    \item \textbf{Sending the BSM result on a quantum wire.} The result is two classical bits.
    \item \textbf{Leaving switch qubits entangled after measurement.} They are consumed and freed by the BSM.
    \item \textbf{Calling swapping teleportation.} Swapping creates a longer ebit; teleportation consumes an ebit to move an unknown state.
    \item \textbf{Treating an optical beam splitter as a qubit copier.} Splitting optical power is not cloning an unknown quantum state.
    \item \textbf{Ignoring failed links.} The swap starts only after the required adjacent links are confirmed.
    \item \textbf{Ignoring time.} Stored pairs must remain coherent while other links and control messages complete.
\end{itemize}

\section{Working Checklist}

\begin{itemize}[leftmargin=*]
    \item[\unchecked] Draw a Bell-pair preparation block.
    \item[\unchecked] Draw a BSM block.
    \item[\unchecked] Draw the two classical output bits.
    \item[\unchecked] Add conditional \(X\) and \(Z\) correction blocks.
    \item[\unchecked] Mark all operations by node.
    \item[\unchecked] Mark the point at which each switch qubit is freed.
    \item[\unchecked] Add heralding before the BSM.
    \item[\unchecked] Add a failure path when only one adjacent link succeeds.
    \item[\unchecked] Add a coherence-time boundary around one complete attempt.
    \item[\unchecked] Explain the circuit from left to right without saying ``the switch sends the qubit.''
\end{itemize}

\section{Circuit-to-Network Translation}

\begin{center}
\begin{tabular}{|p{0.29\textwidth}|p{0.61\textwidth}|}
\hline
\textbf{Circuit element} & \textbf{Network meaning} \\
\hline
Bell-pair preparation & Establish an elementary entangled resource \\
\hline
Qubit wire & Stored quantum state, not necessarily a traveling photon \\
\hline
Measurement & Consumes the measured qubit and creates a classical outcome \\
\hline
Classical wire & Herald, link state, or BSM result \\
\hline
Conditional \(X/Z\) & Endpoint correction or Pauli-frame update \\
\hline
Repeated block & Another time-slot attempt \\
\hline
Parallel Bell blocks & Concurrent neighbor-link attempts \\
\hline
BSM block & Repeater or switch entanglement swap \\
\hline
\end{tabular}
\end{center}

\section{End-of-Week Deliverable}

\begin{itemize}[leftmargin=*]
    \item[\unchecked] One clean Bell-pair circuit.
    \item[\unchecked] One clean BSM circuit.
    \item[\unchecked] One complete Alice--Switch--Bob circuit with ownership labels.
    \item[\unchecked] One paragraph explaining the quantum and classical parts.
    \item[\unchecked] One paragraph naming the physical assumptions hidden by the circuit.
\end{itemize}

\end{document}
